% =========================================================================
% บทที่ 1: หน่วย มาตรฐานทางฟิสิกส์ และการวิเคราะห์ความไม่แน่นอนในการวิจัย
% =========================================================================

\chapter{หน่วย มาตรฐานทางฟิสิกส์ และการวิเคราะห์ความไม่แน่นอนในการวิจัย}
\label{ch:chapter01}

\noindent{\large\bfseries\color{physicsblue} การนิยามระบบหน่วยสากลใหม่ตามค่าคงที่มูลฐานและการประเมินความไม่แน่นอนตามมาตรฐาน GUM}

\vspace{0.4cm}

\begin{equationsbox}{สมการความไม่แน่นอนและการส่งผ่านความคลาดเคลื่อน}
\begin{align}
\bar{x} &= \frac{1}{N}\sum_{i=1}^{N} x_i, \quad s = \sqrt{\frac{1}{N-1}\sum_{i=1}^{N}(x_i - \bar{x})^2} \\
u_A(\bar{x}) &= \frac{s}{\sqrt{N}}, \quad u_B = \frac{a}{\sqrt{3}} \text{ (Uniform)}, \quad u_c(y) = \sqrt{\sum_{i=1}^{n}\left(\frac{\partial f}{\partial x_i}\right)^2 u^2(x_i)}
\end{align}
\end{equationsbox}

\section{ทฤษฎีและกรอบแนวคิดสำคัญ}

การวัดในทางฟิสิกส์มิใช่เพียงการอ่านค่าตัวเลขจากเครื่องมือวัด แต่เป็นกระบวนการเปรียบเทียบเชิงปริมาณกับมาตรฐานสากลที่สามารถสอบกลับได้ (Metrological Traceability) ภายใต้การปรับปรุงระบบหน่วยสากล (SI Redefinition) หน่วยฐานทั้งหมดได้รับการนิยามใหม่โดยตรึงค่าคงที่มูลฐานทางฟิสิกส์ เช่น ค่าคงที่พลังค์ $h = 6.626\,070\,15 \times 10^{-34}\text{ J}\cdot\text{s}$, ความเร็วแสง $c = 299\,792\,458\text{ m/s}$, และประจุธาตุ $e = 1.602\,176\,634 \times 10^{-19}\text{ C}$

ในการรายงานผลการวิจัยทางฟิสิกส์ การระบุค่าความไม่แน่นอนของการวัด (Measurement Uncertainty) ตามคู่มือ ISO/IEC Guide 98-3 (GUM) จำแนกออกเป็นสองประเภทหลัก ได้แก่ ความไม่แน่นอนประเภทเอ (Type A Uncertainty) ซึ่งประเมินด้วยระเบียบวิธีทางสถิติจากการวัดซ้ำ และความไม่แน่นอนประเภทบี (Type B Uncertainty) ซึ่งประเมินจากข้อมูลจำเพาะของเครื่องมือวัด ใบรับรองการสอบเทียบ หรือการแจกแจงความน่าจะเป็นเชิงทฤษฎี

\section{ตัวอย่างการคำนวณตามกรอบวิเคราะห์ P-S-C-T}

\begin{psctexample}{1.1}{การประเมินความไม่แน่นอนรวมของการวัดความหนาแน่นของแท่งโลหะทรงกระบอก}
\textbf{1. ภาพและสถานการณ์ทางกายภาพ (Picture):}\\\\ แท่งโลหะทรงกระบอกมวล $m$ มีเส้นผ่านศูนย์กลาง $d$ และความยาว $L$ ผ่านการวัดละเอียดด้วยไมโครมิเตอร์และเครื่องชั่งดิจิทัล

\vspace{4pt}
\textbf{2. กลยุทธ์และแบบจำลองทางฟิสิกส์ (Strategy):}\\\\ ใช้แบบจำลองทางฟิสิกส์ $\rho = \frac{4m}{\pi d^2 L}$ และใช้กฎการส่งผ่านความไม่แน่นอน (Law of Propagation of Uncertainty) เพื่อคำนวณ $u_c(\rho)$

\vspace{4pt}
\textbf{3. ขั้นตอนการคำนวณเชิงตัวเลข (Calculation):}\\\\ จากฟังก์ชัน $\rho(m, d, L)$ สัมประสิทธิ์ความไว (Sensitivity Coefficients) คำนวณจากอนุพันธ์ย่อย:
\begin{equation}
\frac{\partial \rho}{\partial m} = \frac{4}{\pi d^2 L}, \quad \frac{\partial \rho}{\partial d} = -\frac{8m}{\pi d^3 L}, \quad \frac{\partial \rho}{\partial L} = -\frac{4m}{\pi d^2 L^2}
\end{equation}
ความไม่แน่นอนสัมพัทธ์รวมจึงอยู่ในรูป:
\begin{equation}
\left[\frac{u_c(\rho)}{\rho}\right]^2 = \left[\frac{u(m)}{m}\right]^2 + 4\left[\frac{u(d)}{d}\right]^2 + \left[\frac{u(L)}{L}\right]^2
\end{equation}
สมมติค่าที่วัดได้ $m = 150.00 \pm 0.05\text{ g}$, $d = 20.00 \pm 0.01\text{ mm}$, $L = 60.00 \pm 0.02\text{ mm}$:
\begin{equation}
\frac{u_c(\rho)}{\rho} = \sqrt{\left(\frac{0.05}{150}\right)^2 + 4\left(\frac{0.01}{20}\right)^2 + \left(\frac{0.02}{60}\right)^2} \approx 0.001\,1
\end{equation}
ความหนาแน่น $\rho = 7.9577\text{ g/cm}^3$ และ $u_c(\rho) = 0.0088\text{ g/cm}^3$ บันทึกผลเป็น $\rho = 7.958 \pm 0.018\text{ g/cm}^3$ ที่ระดับความเชื่อมั่น 95\% ($k=2$)

\vspace{4pt}
\textbf{4. การทดสอบความสมเหตุสมผล (Test):}\\\\ มิติของความหนาแน่นตรงตาม $\text{M}\cdot\text{L}^{-3}$ และความไม่แน่นอนสอดคล้องกับพิกัดความคลาดเคลื่อนของเครื่องมือวัดละเอียดระดับไมครอน
\end{psctexample}

\section{การเชื่อมโยงสู่การวิจัยฟิสิกส์ร่วมสมัย}

\begin{researchbox}{ข้อมูลเชิงลึกจากงานวิจัย}
ในการทำวิจัยด้านดาราศาสตร์ฟิสิกส์ด้วยเทคนิค CCD Photometry (เช่น การศึกษาระบบดาวคู่สัมผัส DF Hydrae ของ ผศ.ดร.ชีวะ ทัศนา) การวัดฟลักซ์ความสว่างของดาวฤกษ์ต้องประเมินทั้งสัญญาณรบกวนโฟตอน (Poisson Photon Noise), สัญญาณรบกวนกระแสมืด (Dark Current), และสัญญาณรบกวนจากการอ่านข้อมูล (Readout Noise) ภายใต้แบบจำลอง CCD Equation ซึ่งเทียบเคียงได้กับกฎการส่งผ่านความไม่แน่นอนตามมาตรฐาน GUM
\end{researchbox}

\section{การฝึกแก้โจทย์ปัญหาเชิงระบบตามกรอบ C-C-A-F}

\begin{ccafsolution}{การวิเคราะห์ความคลาดเคลื่อนในการวัดค่าสนามโน้มถ่วงด้วยลูกตุ้มฟิสิกส์}
\textbf{ขั้นที่ 1: การสร้างมโนทัศน์ (Conceptualize):}\\\\ พิจารณาลูกตุ้มอย่างง่ายที่มีความยาว $L$ คาบการแกว่ง $T$ สัมพันธ์กับค่า $g$ ผ่านสมการ $T = 2\pi\sqrt{L/g}$

\vspace{4pt}
\textbf{ขั้นที่ 2: การจำแนกประเภทโจทย์ (Categorize):}\\\\ ปัญหาการส่งผ่านความคลาดเคลื่อนในการวัดทางอ้อม (Indirect Measurement with Error Propagation)

\vspace{4pt}
\textbf{ขั้นที่ 3: การวิเคราะห์เชิงคณิตศาสตร์ (Analyze):}\\\\ จัดรูปสมการหาค่าความเร่งโน้มถ่วง:
\begin{equation}
g = \frac{4\pi^2 L}{T^2} \implies \frac{\Delta g}{g} = \sqrt{\left(\frac{\Delta L}{L}\right)^2 + 4\left(\frac{\Delta T}{T}\right)^2}
\end{equation}
จะเห็นว่าความคลาดเคลื่อนในการจับเวลาคาบ $T$ ส่งผลกระทบต่อความแม่นยำของ $g$ เป็นสองเท่าของความยาว $L$ ดังนั้น ในการออกแบบการทดลอง ผู้วิจัยจึงต้องวัดเวลาการแกว่งต่อเนื่อง 20--50 รอบ เพื่อลดค่า $u_A(T)$

\vspace{4pt}
\textbf{ขั้นที่ 4: การสรุปและประเมินผล (Finalize):}\\\\ การเพิ่มจำนวนรอบการแกว่งเป็นกลยุทธ์สำคัญทางระเบียบวิธีวิจัยที่ช่วยลดความไม่แน่นอนของการวัดได้ตามทฤษฎี $1/\sqrt{N}$
\end{ccafsolution}

\section{แบบฝึกหัดทบทวนและโจทย์ท้าทายประจำบท}

\begin{enumerate}[leftmargin=1.5cm, label={\textbf{\thechapter.\arabic*}}, itemsep=8pt]
    \item \textbf{โจทย์เชิงแนวคิด (Conceptual):} จงอธิบายความสำคัญทางกายภาพของกฎและสมการหลักในบทเรียนนี้ พร้อมยกตัวอย่างสถานการณ์จริงในธรรมชาติ
    \item \textbf{โจทย์การประมาณค่าเชิงฟิสิกส์ (Estimation):} จงใช้การวิเคราะห์เชิงมิติและอันดับของขนาด (Order of Magnitude) ในการประมาณค่าปริมาณสำคัญที่เกี่ยวข้อง
    \item \textbf{โจทย์วิจัยขั้นสูง (Research Challenge):} จงเขียนสคริปต์ Python จำลองพฤติกรรมของระบบตามสมการที่ได้ศึกษา พร้อมพลอตกราฟเปรียบเทียบผลลัพธ์เชิงทฤษฎี
\end{enumerate}

